\documentclass[11pt]{article}
\usepackage{amsmath,amssymb,amsthm}
\newtheorem{theorem}{Theorem}
\title{A joint report deadline for a single-file formation}
\begin{document}
\maketitle

Three possible service requirements can each be met in pairs at the
instant a report arrives, with a standby chosen for the pair.
A single formation for all three instead needs time to move. The
joint obstruction is a cycle of constraints running across the
possible continuations, not a conflict within a single one.

\paragraph{Fixed laws.}
There are \(n\) scalar positions. The allowed formation region \(C\)
is a nonempty compact set specified by difference inequalities
\(x_v-x_u\le b_{uv}\), including individual bounds using an extra
coordinate fixed to zero. These can express clearance, maintaining
a geometric link, or staying in a passage. A truthful report at time
zero selects one nonempty closed terminal set \(F_j\subseteq C\) from
a finite nonempty family. The formation may stand at any point of
\(C\) before the report, or proceed with any shared feasible motion
that does not distinguish the possible reports. The law from the
report until the deadline \(T\) is a double integrator per coordinate.
Velocity is absolutely continuous, with bounds \(a_i=\kappa_i a\)
and \(V_i=\kappa_i V\), for \(a,V,\kappa_i>0\), and everyone must
stop at a point of \(F_j\) by \(T\) while remaining in \(C\). All
robots receive the report at the cut. This does not model staggered
notice, actual radio service or failed actuators.

\begin{theorem}
Let \(D(T)=\int_0^T\min(at,a(T-t),V)\,dt\), the double-integrator
rest-to-rest distance with the base bounds. The joint deadline is
feasible iff there are \(x\in C\) and \(y_j\in F_j\) satisfying
\[
                  |x_i-y_{j,i}|\le\kappa_i D(T)\quad\text{for all }i,j.
\]
This holds even if the pre-report phase can choose a nonzero shared
velocity at the report.
\end{theorem}
\begin{proof}
For any common pre-report state \((z,w)\) and feasible branch,
compare its coordinate velocity to \(w_i\) up to
\(\tau=\min(T/2,V/a)\) and to zero afterward. The respective
differences in absolute value are bounded by
\(\kappa_i\min(at,a(T-t),V)\), using the initial and terminal
velocities and the speed cap. Hence all terminal points belong
to one axis-aligned box centered at \(z+\tau w\) with radii
\(r_i=\kappa_i D(T)\). Their range in each coordinate, denoted
\(\Delta_i\), is at most \(2r_i\).

It remains to choose a center in \(C\) for those same terminal
points. Let \(L_i=\max_j y_{j,i}-r_i\) and
\(U_i=\min_j y_{j,i}+r_i\), including \(L_0=U_0=0\) for the
fixed coordinate. Write \(d_{uv}\) for the shortest-path bound
on \(x_v-x_u\) implied by \(C\), with \(d_{uv}=+\infty\) if none
exists. Every terminal point satisfies it. Thus
\(\max_j y_{j,v}-\min_j y_{j,u}-d_{uv}\le\min(\Delta_v,\Delta_u)\):
take the maximum bound for \(v\) versus \(u\), and separately the
minimum bound, to get the two range estimates. It follows that
\(L_v-U_u\le d_{uv}\), since \(2\min(r_v,r_u)\le r_v+r_u\).
The vector \(x_u=\max_{v:d_{uv}<+\infty}(L_v-d_{uv})\) lies
between \(L\) and \(U\) and obeys each original inequality by
the shortest-path triangle inequality. This proves necessity
inside \(C\), not merely in the ambient space.

For sufficiency, start from \(x\) at rest. In branch \(j\) travel
on the segment to \(y_j\) using the same normalized progress for
every coordinate: at each moment its rate is
\(\min(at,a(T-t),V)/D(T)\). Coordinate bounds follow from the
displacement constraints; convexity keeps the whole segment in
\(C\). All velocities are zero at the endpoints. If \(T=0\),
the inequalities simply require a common terminal point.
\end{proof}

\paragraph{Deadline from a negative cycle.}
When every \(F_j\) is also specified by difference inequalities,
build one block for the standby and one for each terminal. The
travel bounds are edges with weights \(\kappa_i r\), where
\(r=D(T)\). All other weights are fixed. Feasibility is exactly
absence of a negative cycle. A cycle with weight \(c+kr\) requires
\(r\ge -c/k\) if \(k>0\); a negative fixed cycle is infeasible at
all deadlines. The replay raises \(r\) to the lower bound from each
witnessing cycle until the distance labels give a feasible assignment.
Only finitely many simple cycle bounds can be encountered. The
algorithm gives exact displacement certificates; no polynomial
iteration bound is claimed for this update rule.

Fix three robots in \([0,12]\) with both gaps in \([1,3]\) throughout.
The possible terminal laws require, respectively, first gap \(3\),
second gap \(3\), or total span at most \(4\). Every pair can use
one static position. Across all three, the same standby's first
and second gaps must each be at least \(3-2r\), and its total span
at most \(4+2r\). The three laws together force \(r\ge1/3\).
At equality, both standby gaps are \(7/3\); moving the adjacent
pairs outward in the first two worlds, or the outer pair inward
in the third, gives the replay's exact terminal points. With
\(a=1,V=1/2,\kappa_i=1\), the minimum time is \(7/6\). The proof
permits all causal pre-report motions under the stated interface.
The linear-program baseline with the same endpoints and distance
law agrees on 180 fixture variations. The continuous path guarantee
comes from the proof, not the numerical LP.

\paragraph{Why the actuator shape restriction matters.}
If two robots must keep one exact gap, their velocities and
accelerations must coincide. With bounds \((a_1,V_1)=(1,100)\)
and \((a_2,V_2)=(100,1)\), each can advance \(7/2\) alone in time
4 from rest to rest. Together the effective bounds are \((1,1)\)
and they can advance at most 3. The common-box necessity alone
does not prove a compatible path for different speed-to-acceleration
ratios. The present solver makes no claim there. The
result is about a deadline under joint constraints, not a
performance comparison with a control barrier function
implementation or an independent field trial.
\end{document}
